\documentclass[12pt]{article}
\usepackage{amsmath, amssymb}
\usepackage[margin=1in]{geometry}
\usepackage{parskip}

\title{ORF 309 -- Homework 3, Q5a}
\author{}
\date{}

\begin{document}
\maketitle

\section*{Background: Expected Value}

For a discrete random variable $X$, the expected value is a probability-weighted average over all possible values:
\[
    E[X] = \sum_x x \cdot P(X = x)
\]
For example, for a fair die with outcomes $\{1,2,3,4,5,6\}$ each with probability $\frac{1}{6}$:
\[
    E[X] = 1\cdot\tfrac{1}{6} + 2\cdot\tfrac{1}{6} + 3\cdot\tfrac{1}{6} + 4\cdot\tfrac{1}{6} + 5\cdot\tfrac{1}{6} + 6\cdot\tfrac{1}{6} = 3.5
\]
For $X \sim \text{Binomial}(n, p)$, there is a closed-form shortcut:
\[
    E[X] = np
\]

\section*{Q5(a): Probability the Plane is Overbooked}

\textbf{Setup.} 202 tickets are sold. Each passenger independently fails to show up with probability 0.02, so each passenger shows up with probability $p = 0.98$. Let
\[
    X = \text{number of passengers who show up.}
\]
Then $X \sim \text{Binomial}(202,\, 0.98)$.

The plane is overbooked when more than 200 passengers show up, i.e., $X > 200$. Since $X \in \{0, 1, \ldots, 202\}$, this means $X = 201$ or $X = 202$.

\textbf{Calculation.}
\begin{align*}
    P(\text{overbooked}) &= P(X = 201) + P(X = 202) \\[6pt]
    P(X = 201) &= \binom{202}{201}(0.98)^{201}(0.02)^{1} = 202 \cdot (0.98)^{201} \cdot 0.02 \\[6pt]
    P(X = 202) &= \binom{202}{202}(0.98)^{202}(0.02)^{0} = (0.98)^{202}
\end{align*}

\textbf{Equivalent formulation.} Letting $Y = 202 - X$ be the number of no-shows, we have $Y \sim \text{Binomial}(202,\, 0.02)$. Overbooking occurs when $Y < 2$:
\begin{align*}
    P(\text{overbooked}) &= P(Y = 0) + P(Y = 1) \\
    &= (0.98)^{202} + 202 \cdot (0.02) \cdot (0.98)^{201}
\end{align*}

\textbf{Numerical answer.}
\begin{align*}
    (0.98)^{201} &\approx 0.017236 \\
    P(X = 202) &= (0.98)^{202} \approx 0.016891 \\
    P(X = 201) &= 202 \cdot 0.02 \cdot 0.017236 \approx 0.069634
\end{align*}
\[
    \boxed{P(\text{overbooked}) \approx 0.016891 + 0.069634 \approx 0.0865}
\]

\section*{Q5(b): Expected Ticket Revenue}

The airline collects \$500 per \textit{seated} passenger, and the plane seats at most 200. So the number of seated passengers is $\min(X, 200)$ and:
\[
    E[\text{revenue}] = 500 \cdot E[\min(X, 200)]
\]

\textbf{Computing $E[\min(X, 200)]$.} Use the identity $\min(X, 200) = X - \max(X - 200,\, 0)$:
\[
    E[\min(X, 200)] = E[X] - E[\max(X - 200,\, 0)]
\]

\textbf{First term:} By the Binomial shortcut,
\[
    E[X] = np = 202 \times 0.98 = 197.96
\]

\textbf{Second term:} $\max(X-200, 0)$ is nonzero only when $X > 200$. Since $X \leq 202$:
\begin{align*}
    E[\max(X-200,\, 0)] &= 1 \cdot P(X = 201) + 2 \cdot P(X = 202) \\
    &\approx 1(0.069634) + 2(0.016891) \\
    &\approx 0.103416
\end{align*}

\textbf{Putting it together:}
\begin{align*}
    E[\min(X, 200)] &= 197.96 - 0.103416 \approx 197.8566
\end{align*}
\[
    \boxed{E[\text{revenue}] = 500 \times 197.8566 \approx \$98{,}928.29}
\]

\end{document}
